\begin{table}[!htbp]
\centering
\caption{Decomposition of the ten-year survey-based inflation-risk-premium discrepancy}
\label{tab:irp_decomposition}
\begin{threeparttable}
\small
\begin{tabular*}{\textwidth}{@{\extracolsep{\fill}}lrrrr}
\toprule
 & \multicolumn{2}{c}{All survey dates} & \multicolumn{2}{c}{2004 onward} \\
\cmidrule(lr){2-3} \cmidrule(lr){4-5}
Component & Mean & RMSE & Mean & RMSE \\ \midrule
Model premium $-$ survey proxy & 0.216 & 0.496 & 0.155 & 0.374 \\
\quad Definition and convexity wedge & 0.202 & 0.269 & 0.210 & 0.305 \\
\quad Fitted-observable proxy $-$ survey proxy & 0.014 & 0.449 & -0.055 & 0.343 \\
\qquad Nominal-yield fit contribution & 0.003 & 0.072 & 0.006 & 0.081 \\
\qquad Real-yield fit contribution & -0.000 & 0.285 & -0.033 & 0.169 \\
\qquad CPI-forecast fit contribution & 0.012 & 0.204 & -0.027 & 0.215 \\
\bottomrule
\end{tabular*}
\begin{tablenotes}[flushleft]
\footnotesize
\item Notes: All entries are annualized percentage points. The survey-based expectation-hypothesis proxy is the observed ten-year nominal yield minus the observed ten-year real yield and the ten-year CPI forecast. The fitted-observable proxy replaces those three inputs with their fitted values. The model premium compares log certainty-equivalent values of cumulative inflation under the risk-neutral and physical measures. Their difference is a definition and convexity wedge generated by Jensen, covariance, and nonlinear lower-bound terms; it is not a fitting error. The fitted-observable error equals exactly the sum of the three fit contributions; their RMSEs do not add because the contributions covary. The columns contain 140 and 91 survey dates, respectively.
\end{tablenotes}
\end{threeparttable}
\end{table}
